\chapter{Troubleshooting}
\label{app:troubleshooting}

Organised by \emph{symptom}, because that is what you have when you arrive here.

% ============================================================
\section{Package and restore}

\subsection{``Package X does not exist''}

Almost always the prerelease flag. A large fraction of the Agent Framework family
carries a \code{-preview}, \code{-alpha} or \code{-rc} suffix even though the core
is GA (Appendix~\ref{app:packages}). NuGet will not resolve those unless you ask
it to.

\begin{shellcmd}
dotnet add package Microsoft.Extensions.AI.Ollama --prerelease
\end{shellcmd}

In Visual Studio: \emph{Include prerelease} in the NuGet Package Manager.

\subsection{``Version conflict'' after adding a hosting or harness package}

You have crossed release trains. Core sits at \mafcore{}; the harness, durable
execution and Foundry hosting sit at \mafext{}. This is expected --- pin both
numbers rather than trying to force one.

If a transitive conflict follows, \code{CentralPackageTransitivePinningEnabled}
in \code{Directory.Packages.props} resolves most of them by pinning the
transitive graph to your declared versions.

\subsection{``NU1008: Projects that use CPM should not define version''}

A \code{PackageReference} still carries a \code{Version} attribute. Central
Package Management forbids it. Strip versions from every \code{.csproj}; they
belong only in \code{Directory.Packages.props}.

\subsection{Build fails only on prerelease packages after enabling \texttt{TreatWarningsAsErrors}}

Preview packages ship analyzers and experimental-API attributes that emit
warnings. Rather than abandoning warnings-as-errors for the solution, suppress
the specific diagnostic in the one project that needs it, and leave a comment
saying which package caused it and when to remove the suppression.

% ============================================================
\section{Authentication}

\subsection{\texttt{DefaultAzureCredential} fails with no obvious cause}

It walks a chain of credential sources and reports the aggregate failure, which
buries the useful part. Narrow it:

\begin{shellcmd}
az login
az account show
\end{shellcmd}

If the CLI session is valid but the credential still fails, force a single source
--- \api{AzureCliCredential} on a dev machine --- so the error becomes specific.
Then check role assignment on the resource: a valid identity with no data-plane
role fails at the first call, not at construction.

\subsection{401 or 403 on the first agent run, not at startup}

Expected. Client construction does not authenticate; the first request does. The
failure surfaces inside \api{RunAsync} rather than where the client was built.

\subsection{Secrets are in \texttt{appsettings.json}}

Move them. \code{dotnet user-secrets init} then \code{dotnet user-secrets set}.
See \S\ref{sec:swap} in Chapter~\ref{ch:first-agent}.

% ============================================================
\section{The agent will not call your tool}
\label{sec:tb-no-tool-call}

The most confusing failure in the framework, because \textbf{nothing errors}. The
agent answers politely and simply never calls the tool. Work through these in
order --- the order matters, because most people start at step four and lose an
afternoon.

\begin{enumerate}[leftmargin=1.4em]
\item \textbf{Does the model support function calling at all?} Many small local
      models do not. Swap to a model documented as supporting tool use ---
      \code{llama3.2} or \code{qwen3:4b} --- and retest before changing anything
      else.
\item \textbf{Is function-invocation middleware wired?} Some clients need
      \api{UseFunctionInvocation} explicitly. See
      Appendix~\ref{app:providers}, Listing~\ref{lst:prov-funcinvoke}.
\item \textbf{Was the tool actually registered?} Break after agent construction
      and inspect the tool collection. A tool that is not there cannot be called.
\item \textbf{Is the description good enough?} Only now. The description is the
      model's entire basis for selection; ``gets data'' will not be selected over
      a plausible-looking alternative. Chapter~\ref{ch:tools} measures how much
      this matters.
\item \textbf{Is the parameter schema too complex?} Deeply nested objects degrade
      selection accuracy sharply on smaller models. Flatten and retest.
\end{enumerate}

% ============================================================
\section{Conversation and state}

\subsection{The agent forgets everything between turns}

You are not passing a thread. \api{RunAsync} without an \api{AgentThread} creates
a throwaway one that lives only for that call. See \S\ref{sec:threads}.

\subsection{Deserialising a thread fails, or behaves oddly}

The thread must be resumed by an \emph{identically configured} agent instance.
If construction differs between the process that saved and the process that
loads --- a different model, different instructions, a different provider --- the
failure may be silent rather than thrown. Register the agent once in DI and build
it from configuration that cannot drift.

\subsection{Serialised threads are far smaller than expected}

The provider is storing conversations server-side, so the serialised thread
contains an identifier rather than the messages. This is correct behaviour, not
data loss --- but it means your backup contains a pointer to state you do not
own, and that the remote thread is yours to delete.

\subsection{Responses degrade during a long tool-calling chain}

Context window pressure. Confirm by logging token counts per turn before
reaching for prompt changes. Chapter~\ref{ch:memory} covers automatic compaction;
be aware that compaction is lossy by design and can remove something you needed.

% ============================================================
\section{Workflows}

\subsection{Type mismatch between executors}

Messages are typed and edges enforce it. Where the compile-time route generators
are in use (\pkg{Microsoft.Agents.AI.Workflows.Generators}) this surfaces at
build time; without them it surfaces at run time, which is worse. Prefer the
generators.

\subsection{Resume produces a different result than an uninterrupted run}

Something outside the checkpoint is carrying state --- a static field, a cached
client, an external store mutated between checkpoint and resume. Checkpointing
captures workflow state, not the state of everything your executors happen to
touch. Make executors take their dependencies explicitly.

\subsection{A human-in-the-loop workflow never resumes}

The response is not reaching the port, or the workflow instance being resumed is
not the one that paused. Log the workflow instance identifier on both sides
before debugging anything else.

% ============================================================
\section{Observability}

\subsection{No traces appear}

In order: is the exporter configured and pointing at a reachable endpoint; is the
activity source registered; is sampling discarding them; is the process exiting
before the exporter flushes? The last one accounts for most console-app cases ---
dispose the tracer provider, or add a delay, before the process ends.

\subsection{Token counts are missing or wrong}

Providers report usage differently, and some do not report it on streaming
responses at all. Verify per provider before building cost calculations on top.
See Appendix~\ref{app:providers}.

\subsection{Traces appear but multi-agent runs are unreadable}

Spans are there but the parent-child relationship is lost, usually because
context is not flowing across an async boundary or a queue. Check that ambient
activity context propagates before concluding the instrumentation is inadequate.

% ============================================================
\section{When the answer is not here}

\begin{itemize}[leftmargin=1.4em]
\item \code{github.com/microsoft/agent-framework/releases} --- read the notes for
      every version between yours and current. Breaking changes are flagged
      \code{[BREAKING]} in the title.
\item \code{github.com/microsoft/agent-framework/discussions}
\item The \code{dotnet/samples} directory --- a working sample is usually faster
      than a documentation page.
\item The framework is open source. When behaviour contradicts documentation, the
      source settles it, and given the release cadence the documentation is
      sometimes the thing that is out of date.
\end{itemize}
